% ECONOMICS_PROBLEM: {"id": "Q16-490", "title": "Agricultural technology adoption constraints", "jel_code": "Q16", "source_id": "OP-0377"}
% !TeX program = xelatex
\documentclass[11pt,a4paper]{article}
\usepackage[margin=23mm]{geometry}
\usepackage{fontspec}
\setmainfont{Latin Modern Roman}
\usepackage{amsmath,amssymb,mathtools}
\usepackage{xcolor,enumitem,booktabs,longtable,array,xurl,hyperref}
\definecolor{muted}{HTML}{53636C}
\hypersetup{colorlinks=true,urlcolor=blue,linkcolor=blue}
\setlength{\parindent}{0pt}
\setlength{\parskip}{5pt}
\newcommand{\R}{\mathbb R}
\newcommand{\E}{\mathbb E}
\newcommand{\N}{\mathbb N}
\newcommand{\ind}{\mathbf 1}
\newcommand{\Var}{\operatorname{Var}}
\newcommand{\Cov}{\operatorname{Cov}}
\newcommand{\supp}{\operatorname{supp}}
\DeclareMathOperator*{\argmin}{arg\,min}
\DeclareMathOperator*{\argmax}{arg\,max}
\newcommand{\entrymeta}[3]{{\sffamily\small\color{muted}\textbf{\texttt{#1}}\quad|\quad #2\par #3\par}}
\newcommand{\Problemref}[1]{\texttt{#1}}
\newcommand{\JELref}[2]{\texttt{#2} (JEL \texttt{#1})}
\begin{document}
\section*{Agricultural technology adoption constraints}
\textbf{Problem ID:} \texttt{Q16-490}\quad \textbf{JEL:} \texttt{Q16}\par
% BEGIN ECONOMICS STATEMENT
\label{problem:OP-0377}
\label{audited:ECON-093}
{\sffamily\small\color{muted}Original subject group: Development and poverty\par}
\label{definitions:problem:30007659}
\textbf{Audit classification:} Published research agenda. Metadata and abstract checked. It expressly flags future directions involving credit, insurance and information; detailed full-text agenda was not inspected and the five-part menu is editorial.
\subsection*{Definitions and precise research target}
\textbf{Complete editorial problem.} Fix smallholder farmers, a crop technology, and several agricultural seasons. Let \(p=(c,r,i,q,m)\) denote a feasible intervention supplying credit, risk coverage, information, verified input quality, and market access. Let \(A_{it}(p)\) indicate technology adoption and \(\Pi_{it}(p)\) net profits after input, labor, finance, and insurance costs. Define \(\tau_\Pi(T,p)=E[\sum_{t=0}^{T}\delta^t\{\Pi_{it}(p)-\Pi_{it}(0)\}]\), where \(\delta\in(0,1]\) is a declared discount factor. Estimate adoption, downside risk, and sustained profitability, allowing interactions between constraints. Adoption itself is not evidence of welfare improvement when returns are heterogeneous. A factorial or sequential randomized design should distinguish learning from financial constraints and insurance demand from uninsurable uncertainty. Specify weather support, input reliability, and crop-price responses before extrapolating across seasons or regions. Determine which cost-effective bundles generate continued use after subsidies and advice stop. Account for peer learning and product-market spillovers through an explicit assignment design. The review identifies constraint interactions and future research, while many individual intervention effects are established. The precise policy menu, objective, and multiseason estimand here are an adapted research specification rather than a source theorem.

\textbf{Definitions and admissible scope.} Define the technology by input, training, and usage requirements, and adoption by verified use on a stated share of land. Net profits deduct all input costs, labor opportunity cost, borrowing, insurance premia, and depreciation; gross revenue is not profit. Risk coverage specifies triggers, basis risk, indemnity, and insurer solvency. Information specifies signal accuracy and delivery. The regime includes subsidy withdrawal dates and market access. Weather and prices have a declared joint stochastic law; all outcomes are integrable. Factorial interactions are contrasts such as \(E[\Pi(1,1)-\Pi(1,0)-\Pi(0,1)+\Pi(0,0)]\), which establish complementarities of interventions under the design but do not automatically name the underlying behavioral mechanism. Take dates \(0,\ldots,T\), \(T<\infty\), a fixed crop-season calendar, baseline farmer weights, and a declared intervention menu. Baseline \(p=0\) denotes no added bundle, retaining existing access. A downside-risk target is, for example, \(\Pr(\sum_t\delta^t\Pi_{it}(p)<0)\), or a stated consumption-shortfall probability; it is separate from \(\tau_\Pi\). Sustained adoption means verified use at every prespecified season after support withdrawal. A welfare criterion additionally specifies risk preferences and consumption opportunities. The two-coordinate factorial shorthand \(\Pi(1,1)\) holds all other bundle coordinates and coverage fixed. Returns and interaction effects are defined only on the joint design support.
\subsection*{Variants and assumption changes}
\begin{enumerate}
\item Credit, insurance, and information (source-motivated): cross the three dimensions in a factorial design and distinguish incomplete learning from insurable production risk and basis risk.
\item Input quality and scale (editorial): add verified-input certification and regional coverage variation. Compare multiseason utility and adoption after support stops, allowing crop prices, peer information, and supplier behavior to change.
\end{enumerate}
\subsection*{References and source audit}
\textbf{Support and access.} Metadata and abstract checked. It expressly flags future directions involving credit, insurance and information; detailed full-text agenda was not inspected and the five-part menu is editorial. \textbf{Locator:} Publisher abstract, final sentence.
\begin{enumerate}\raggedright
\item\hypertarget{definitions:source:30007659:1}{} Jeremy R. Magruder. 2018. \emph{An Assessment of Experimental Evidence on Agricultural Technology Adoption in Developing Countries}. Abstract, paragraphs 1–2.
\url{https://www.annualreviews.org/content/journals/10.1146/annurev-resource-100517-023202}.
DOI: \url{https://doi.org/10.1146/annurev-resource-100517-023202}.
\end{enumerate}

\subsection*{Integrated refinement: ECON-093}
{\sffamily\small\color{muted}Multiple constraints and heterogeneous agricultural technology adoption\par}
This refinement adopts the additional source conventions
(page~\pageref{audited:conventions}), subject to its local restrictions.
\textbf{Heterogeneous welfare gains and bundle interactions.}
For a finite supported menu relaxing credit, risk, knowledge and market access,
identify $\tau_\pi(z,x)=\mathbb E[\pi_i(z)-\pi_i(0)\mid X_i=x]$ with baseline
covariates and an explicit measurement-error model. Fix component versions,
intensity, timing, crop seasons and coverage. Factorial assignment can identify
arm and interaction means under its interference assumptions; naming a
constraint's causal contribution additionally requires credible component
interventions. Compare net profits after finance, labor and insurance costs,
household utility including risk, and social environmental effects separately.
Raising adoption alone need not improve welfare. Report interaction scales and
uncertainty, and allow prices, technology supply and peer learning to change
under widespread adoption.
\subsection*{Supplied source audit for ECON-093}
Local [S1], [S2], etc. in this audit and the references immediately below
belong to ECON-093, independently of the original entry's reference numbering.
[S1] and [S2] directly support multiple constraints and farmer heterogeneity. This conditional bundle-effect formulation is editorial, not a general theorem that relaxing all constraints raises welfare.
\subsection*{References for integrated refinement ECON-093}
{\footnotesize\raggedright
\par\noindent\textbf{[S1]} Tavneet Suri, Chris Udry, Jenny Aker, Chris Barrett, Laura Falcao Bergquist, Michael Carter, Lorenzo Casaburi, Robert Darko Osei, Douglas Gollin, Vivian Hoffmann, Thomas Jayne, Naureen Karachiwalla, Harounan Kazianga, Jeremy Magruder, Hope Michelson, Meredith Startz, and Emilia Tjernstrom (2024). \emph{Agricultural Technology in Africa}. VoxDevLit 5(2), March 2024.\par \textit{Verified locator / source role:} Primary publisher page checked for review identity, authors, version and background. The specific published agenda is corroborated in [S2]; the exact formal target is editorial.\par \url{https://voxdev.org/voxdevlit/agricultural-technology-africa}\par\smallskip
\par\noindent\textbf{[S2]} Oliver Hanney / VoxDev (living compilation). \emph{Evidence gaps: Key unanswered questions in development economics}. Accessed 8 October 2026. The page currently covers 20 topics; the ``16-topics URL is a historical slug.\par \textit{Verified locator / source role:} Topic heading: Agricultural Technology in Africa. Supports the broad research agenda; this is not a verbatim quotation or an attribution of the displayed mathematical target.\par \url{https://voxdev.org/topic/evidence-gaps-key-unanswered-questions-across-16-topics-development-economics}\par\smallskip
}

\subsection*{Common conventions: Source mathematical conventions}

These conventions apply to each entry unless explicitly replaced.

\textbf{Models and identification.} A primitive model \(m\) specifies all probability laws, feasible choices, information, equilibrium and measurement rules used by its target. The admissible class \(\mathcal M\) is fixed independently of outcomes. If \(P_O(m)\) is its observable probability law and \(\tau(m)\) the target, its identified set at observable law \(P\) is
\[
 \mathcal I_\tau(P)=\{\tau(m):m\in\mathcal M,\ P_O(m)=P\}.
\]
Point identification means that this set is a singleton. An empty set signals incompatibility of data and assumptions. Coordinate bounds need not describe the joint set. Bounds are sharp only if every retained target is attainable or arbitrarily approachable by compatible models. Finite bounds require explicit restrictions on unobserved quantities; unrestricted confounding, hidden wealth, extrapolation or arbitrary equilibrium selection can yield uninformative sets.

\textbf{Causal interventions.} A potential outcome \(Y_i(p)\) is the outcome under a complete policy regime \(p\), including timing, financing, information and market spillovers. The target population and units are fixed before treatment. With interference, use the complete assignment vector or a justified exposure mapping; individual no-interference notation alone cannot identify an economy-wide intervention. A difference of mean potential outcomes does not require the cross-world joint distribution. Conditioning on post-treatment migration, survival, enrollment or employment changes the estimand and needs separate justification.

\textbf{Statistical statements.} An estimator, confidence set, forecast or test specifies a sampling law, model class and loss/coverage criterion. Uniform \((1-\alpha)\)-coverage means \(\inf_{m\in\mathcal M}\Pr_m\{\tau(m)\in C_n\}\geq1-\alpha\), or an explicitly asymptotic version. The whole parameter space trivially has coverage, so informativeness requires a stated width, risk or power objective. A forecasting improvement is relative to a named benchmark, information set and evaluation population.

\textbf{Equilibrium and optimization.} An equilibrium comprises feasible plans, optimality, beliefs where needed, budget constraints and all market/resource clearing. The primitives or a stated benchmark define each correspondence \(\mathcal E(p;m)\). Nonemptiness is assumed or proved, never inferred just from compact policy sets. A selected-equilibrium target uses a declared selection rule; a robust target quantifies over all permitted equilibria. Use suprema or infima unless attainment is established. An \(\varepsilon\)-optimal policy is feasible and within \(\varepsilon\) of the finite optimum value. Report an empty feasible set separately.

\textbf{Welfare and accounting.} Normative utility, interpersonal normalization, social weights, numeraire, discounting and the use of public receipts are primitives. Behavioral utility need not coincide with the welfare criterion. Household welfare includes ownership income and rents once; adding profits again to compensating variation generally double-counts them. A monetary equivalent is defined by an explicit transfer and price convention, with monotonicity and range conditions. Average effects, mechanism contributions, transfers and welfare losses are different targets. Infinite sums require convergence and dynamic feasible plans require terminal or no-Ponzi conditions.

\textbf{Source versions.} A working-paper date, revision date and journal publication year are distinguished. A DOI for a journal version is not automatically the DOI of an earlier manuscript. Locators refer to the stated version and printed pages where specified. A metadata-only check does not verify a claim about a full-text conclusion. Every entry's source subsection narrows the provenance claim accordingly.

\subsection*{Common conventions: Additional-source status and mathematical conventions}

\label{audited:conventions}

Of the 100 supplied ECON entries, 65 distinct problems appear here; 32 close
matches refine earlier entries, and three duplicate questions map to existing
entries. Appendix~\ref{app:audited-crosswalk} lists every destination.
The attachment's P/R/S/U distinctions and source audits are retained. Its
precise mathematical formalizations are editorial unless the individual audit
says otherwise. Candidate subsidy targets ECON-004--006 retain their U flags;
the resolved approval-core question ECON-007 updates OP-0202 above.
The following supplied conventions govern both this part and the attributed
ECON refinements elsewhere in the catalogue, subject to local restrictions.

\textbf{Parameterized questions.} A problem family lists inputs that must be fixed before an instance is evaluated: population, horizon, policies, information, data law, model class and welfare weights. These are required choices, not quantities the citations are assumed to specify. Undefined applications should not be treated as identified estimands.

\textbf{Indices and integrability.} Sets of agents, firms and regions are finite unless a population measure is explicitly used. \(i,j\) index units, \(r\) regions, \(t\) dates and \(h\) horizons. \(\mathbb E\) and \(\Pr\) refer to the declared population and law; \(\mathbf1\{A\}\) is an indicator. Every displayed expectation and discounted sum needs existence in the stated domain. Infinite horizons use a discount in \((0,1)\) and a convergence condition; finite horizons may use discount one.

\textbf{Optimization and existence.} A displayed \(\max\), \(\min\), or \(\arg\max\) is conditional on attainment. For finite feasible menus this is immediate; general spaces need continuity and compactness/coercivity or another existence argument. Without attainment, the target is the supremum/infimum and arbitrarily near-optimal feasible choices. \(\inf\varnothing=+\infty\). Empty feasibility and an unidentified target are different issues.

\textbf{Potential outcomes.} \(Y_i(z)\) is the outcome under a specified intervention, not the observed conditional mean among people choosing \(z\). In binary comparisons \(\Delta Y_i=Y_i(1)-Y_i(0)\). Specify assignment, treatment versions, compliance, information and observation horizon. Interference is allowed under a full policy vector; compressed exposure notation needs a justified mapping. No interference, consistency and overlap are substantive assumptions, not automatic consequences of notation.

\textbf{Populations and observation.} Fix baseline eligibility and population weights. Conditioning on post-policy employment, survival, relocation or program uptake can change the population being compared. Death is not ordinary loss to follow-up; outcomes truncated by death need a composite convention, joint survival target, or explicitly defined survivor stratum. Fertility-changing interventions need a population functional rather than an unexplained fixed descendant cohort. Measurement and missingness models must be supplied.

\textbf{Identification.} A structural model \(m\in\mathcal M\) implies observable law \(P_m\) and target \(\theta(m)\). Identification means
\[
 P_{m_1}=P_{m_2}\ \Longrightarrow\ \theta(m_1)=\theta(m_2)
 \quad\text{for all }m_1,m_2\in\mathcal M .
\]
Without identification, the target is
\[
 \Theta(P;\mathcal M)=\{\theta(m):m\in\mathcal M,\ P_m=P\}.
\]
Writing \(\theta(P)\) before establishing identification can hide incompatible latent models. Confidence sets additionally need a sampling law, confidence level, asymptotic sequence and uniformity class. Point identification, coverage and informative precision are separate properties.

\textbf{Mediation.} Total effects of randomized assignment are distinct from cross-world natural indirect effects. Belief/effort-channel effects require a meaningful mediator intervention and additional measurement, exclusion or mediation assumptions. Randomization of the initial policy alone does not supply those assumptions. Report bounds or interventional alternatives when the stronger target is unsupported.

\textbf{Equilibrium.} A counterfactual \(W(z)\), \(p(z)\), or \(\mu_t^z\) requires individual optimization, feasibility, government/resource budgets, market clearing and state-transition laws. State equilibrium existence and selection, or report ranges over compatible equilibria. Initial-state integration includes subsequent endogenous transitions; current-state integration must use the policy-dependent distribution. Reduced-form invariance after policy is not assumed.

\textbf{Welfare accounts.} Scores, utility, health, dollars and ecological quantities require explicit conversion before addition. Fees, taxes, subsidies, wages and extortion payments are transfers between parties; distributional weights can make them welfare relevant, but their face value is not automatically a resource loss. State ownership, geographic boundaries, foreign-income weights and fiscal finance. Do not add profits to household welfare already including dividends, or count land capitalization and the underlying benefit twice. A benefit-cost expression must distinguish gross benefits from net welfare and subtract every real cost exactly once.

\textbf{Derivatives and risk.} Log derivatives require positive arguments; local responses require admissible perturbations and specified equilibrium branches. Differentiation under expectations needs regularity/domination, and boundaries may allow only one-sided derivatives. Risk uses a specified law; ambiguity uses a declared nonempty law set on a common history space. Adaptive policies must be nonanticipative. A robust criterion is an editorial choice unless attributed explicitly.

\textbf{Scope overlaps.} Allocation, collective choice, contracts, household bargaining and coordinated policy overlap economics with optimization and game theory. They are retained because they appeared in the original catalogue; no claim of a strictly game-theory-free selection is made.
% END ECONOMICS STATEMENT
\end{document}
